\chapter{Thảo luận, kết luận và hướng phát triển}
\label{chap:discussion}

Các thí nghiệm trong Chương~\ref{chap:results} trả lời một câu hỏi thiết kế cụ thể: có thể khai thác thông tin đa phương thức bằng một interaction graph gọn hơn mà vẫn giữ hiệu năng validation cạnh tranh hay không? Chương này diễn giải các bằng chứng, phân biệt kết luận được hỗ trợ với giả thuyết cần kiểm chứng thêm, và xác định phạm vi sử dụng của mô hình.

\section{Thảo luận}
Kết quả cho thấy dùng ba MRI chưa mang lại lợi ích ổn định so với một MRI. Ghép thêm radiomics và dữ liệu lâm sàng tạo tín hiệu tích cực trên validation, nhưng các giai đoạn còn khác cấu hình huấn luyện nên chưa thể tách riêng lợi ích của từng thành phần. Việc thêm thời điểm hay nguồn thông tin cần được kiểm tra trong cùng điều kiện thí nghiệm.

CP phân rã tensor trọng số của lớp tương tác đầu tiên \cite{kolda2009tensor}, giúp giảm mạnh số tham số của khối này. Trong ba seed, điểm validation trung bình giảm nhẹ so với full outer-product; kết quả chưa cho thấy CP tự cải thiện dự báo. Vì MRI encoder vẫn chiếm phần lớn tham số, mức giảm của toàn mạng nhỏ hơn nhiều. Lợi ích về tốc độ và bộ nhớ cần được đo trực tiếp.

RC-Free được chọn sau ablation một seed và kiểm tra lại bằng ba seed. Mô hình giữ đủ MRI, radiomics và dữ liệu lâm sàng, chỉ bỏ tương tác trực tiếp radiomics--clinical. MRI kết nối với hai nhánh còn lại qua MR và MC (Hình~\ref{fig:mri-hub}); gate điều chỉnh đóng góp theo đầu vào, còn residual giữ thông tin gốc. Kết quả ủng hộ một cấu trúc gọn có hiệu năng validation cạnh tranh trên bộ dữ liệu hiện tại, chưa chứng minh RC vô ích trên mọi dữ liệu. Gate và attention cũng chưa được kiểm chứng như chỉ dấu lâm sàng.

\begin{figure}[htbp]
\centering
\includegraphics[width=\textwidth]{mri-hub-topology.pdf}
\caption[MRI làm trung tâm của topology RC-Free]{Topology cuối được lựa chọn qua validation và ablation: mô hình giữ đủ MRI, radiomics và clinical, sử dụng MRI làm trung tâm tương tác sau khi bỏ direct radiomics--clinical pair. Các cạnh biểu diễn trao đổi feature qua MR/MC, không phải quan hệ nhân quả hay kết nối giải phẫu.}
\label{fig:mri-hub}
\end{figure}

Cox elastic-net là đối chứng mạnh cho dữ liệu bảng. Tuy nhiên, khác biệt về cách tính C-index và chọn mô hình khiến so sánh với mạng sâu chưa hoàn toàn đồng nhất. Vì vậy, các kết luận chỉ áp dụng trong điều kiện thí nghiệm hiện có.

\section{Giới hạn}
\begin{enumerate}
\item \textbf{Dữ liệu và thời điểm dự báo.} Mẫu nghiên cứu gồm 359 người ADNI có đủ ba MRI trong 18 tháng và không tiến triển sớm. Biến cố dựa trên ngưỡng MMSE/CDR, chưa thay thế chẩn đoán AD được xác nhận. Mốc dự báo là MRI thứ ba, nhưng điều kiện chọn mẫu được kiểm tra đến hết 18 tháng; dữ liệu lâm sàng ghép trong $\pm90$ ngày còn có thể xuất hiện sau MRI. Vì vậy, kết quả chưa chứng minh khả năng dự báo bằng dữ liệu có sẵn đúng tại thời điểm đó.
\item \textbf{Độ chắc chắn của kết quả.} Tập validation chỉ có 54 người và 16 biến cố. Ba seed dùng cùng một phép chia dữ liệu; độ lệch chuẩn giữa seed không phải khoảng tin cậy. Ablation ban đầu chỉ chạy một seed, còn nhiều lần chọn mô hình trên validation có thể làm hiệu năng được đánh giá quá cao. Một số đối chứng khác cấu hình huấn luyện; Cox loss tính trên từng nhóm mẫu (batch) và C-index chỉ đánh giá thứ hạng nguy cơ, chưa đánh giá độ chính xác của xác suất dự báo.
\item \textbf{Kiểm chứng trên dữ liệu mới và hiệu quả triển khai.} Các thí nghiệm trước từng dùng tập test; đánh giá cuối cùng theo quy trình chính thức vẫn chưa được thực hiện. Mô hình chưa được kiểm tra trên nhóm người độc lập, khi thiếu một nguồn dữ liệu hay khi chất lượng MRI thay đổi. Rank CP còn cố định, chưa so sánh nhiều giá trị trong cùng điều kiện huấn luyện. Mức giảm tham số cũng chưa chứng minh mức tiết kiệm thời gian hay bộ nhớ vì chưa có phép đo đồng nhất.
\end{enumerate}

\section{Kết luận}
Khóa luận phát triển một pipeline survival từ MRI-only đến longitudinal MRI, ghép nối đa phương thức, dynamic pairwise fusion, CP-factorized interaction và ablation topology. Quá trình này xuất phát từ câu hỏi của nghiên cứu longitudinal MRI--radiomics: bổ sung một modality bằng ghép nối chưa bảo đảm khai thác được tín hiệu của nó \cite{aghajanian2025longitudinal}.

Kiến trúc được chọn là RC-Free CP-Factorized Pairwise Fusion. Mô hình giữ MRI, radiomics và clinical, tham số hóa MR/MC bằng CP, đưa tương tác trở lại modality qua gated residual, rồi xử lý ba visit bằng T-LSTM và attention trước Cox head. Bằng chứng validation hiện tại ủng hộ một graph gọn có hiệu năng cạnh tranh; CP giảm chi phí của interaction core, còn loại RC đơn giản hóa topology. Các kết luận này là kết luận về thiết kế trong cohort hiện tại, chưa phải xác nhận ưu thế lâm sàng hoặc kết quả test.

\section{Hướng phát triển}
Ưu tiên đầu tiên là xây dựng và đánh giá protocol tiến cứu với time origin rõ ràng, eligibility phù hợp lúc dự báo và clinical chỉ lấy từ các ngày đã quan sát; xác nhận endpoint bằng chẩn đoán có thẩm định nếu có. Cần giữ provenance của dữ liệu hiện tại và tạo phiên bản cohort mới thay vì sửa âm thầm bản đóng băng.

Sau khi khóa protocol, nên thực hiện đối chứng MRI/temporal/modality trên cấu hình đồng nhất, mở rộng số split hoặc cohort độc lập, và chạy formal final evaluation có công bố lịch sử test. Calibration, time-dependent Brier score và các metric theo horizon sẽ bổ sung cho concordance \cite{graf1999brier}. Khoảng tin cậy theo người bệnh cần được tách khỏi SD theo seed.

Các hướng về kiến trúc gồm xử lý thiếu modality, rank-adaptive interaction, uncertainty-aware routing và temporal encoders mới khi quy mô dữ liệu cho phép. Pretrained brain representations như hướng brain latent diffusion có thể hữu ích \cite{li2026longitudinal}, nhưng mọi cải tiến cần được kiểm tra bằng đối chứng phù hợp. Cuối cùng, đo latency, memory và độ ổn định trên phần cứng cố định sẽ làm rõ lợi ích triển khai thực sự của kiến trúc gọn.
